\begin{center}
\textsc{Gauss} to \textsc{Bessel}.\quad Göttingen, 20~November 1826.
\end{center}

\par\noindent\hspace*{\parindent}.\ .\ .\ .\ .\ .\quad Working through the material for the projected work on my measurements takes up much of my time.
The calculations for my principal triangles, comprising 33 points, were in fact finished long ago, but much work is needed to compute the many secondary points located by intersection, which are important for the geography of a considerable part of northern Germany, since they have so far either not been computed at all or have been computed only provisionally.
My list now contains about 250 points, all from the more northerly part of my measurements.
Since my triangles are connected through those of Electoral Hesse to those of Bavaria and Württemberg, and \textsc{Bohnenberger} has also kindly sent me the latter (as \textsc{Eckhardt} has those of Darmstadt), I am very much vexed that all my efforts to obtain the Bavarian ones
\par\smallskip\noindent{\footnotesize IX.}\hfill{\footnotesize46}\par
\clearpage
\noindent362\hfill{\small\textbf{CORRESPONDENCE WITH BESSEL.}}\par\medskip
\noindent have so far been in vain. \textsc{Schumacher} told me that \textsc{Soldner} had said to him that the reason the commission in Munich had not granted my request to have them sent to me was that it supposed that I wanted to do calculations on these triangles!
I have had the same experience with the Austrian ones.
General \textsc{von Müffling} ought also to make public his triangles east of Seeberg, which connect Berlin and extend across Silesia.

\par\noindent\hspace*{\parindent}The far more extensive theoretical part causes me still greater difficulty, as it reaches into other parts of mathematics in so many ways.
I see no other remedy here than to detach several large principal portions from the work, so that they can be developed independently and in proper detail.
In a sense I have already made a start on this in my paper on the mapping of surfaces while preserving similarity of the smallest parts. A second treatise, which I presented to the royal society a few months ago and which will, I hope, soon be printed, contains the principles and methods for adjusting the measurements (incidentally, one can also see indirectly from it how far \textsc{Krayenhoff}'s measurements fall short of the accuracy wrongly attributed to them).
Perhaps I will first prepare a third treatise, which will develop various new theorems on curved surfaces, shortest lines, the representation of curved surfaces in the plane, etc.
If all these subjects had had to be included in my projected work, I would either have had to dispatch some things without treating them thoroughly or to give the work a very motley appearance.\ .\ .\ .\ .\ .\ .
